\documentclass{standalone}
\usepackage{graphicx}	
\usepackage{amssymb, amsmath, amsthm}
\usepackage{color}

\usepackage{tikz}
\usetikzlibrary{arrows.meta, math}

\definecolor{light}{RGB}{220, 188, 188}
\definecolor{mid}{RGB}{185, 124, 124}
\definecolor{dark}{RGB}{143, 39, 39}
\definecolor{highlight}{RGB}{180, 31, 180}
\definecolor{lightteal}{RGB}{107, 142, 142}
\definecolor{midteal}{RGB}{72, 117, 117}
\definecolor{darkteal}{RGB}{29, 79, 79}
\definecolor{gray10}{gray}{0.1}
\definecolor{gray20}{gray}{0.2}
\definecolor{gray30}{gray}{0.3}
\definecolor{gray40}{gray}{0.4}
\definecolor{gray60}{gray}{0.6}
\definecolor{gray70}{gray}{0.7}
\definecolor{gray80}{gray}{0.8}
\definecolor{gray90}{gray}{0.9}
\definecolor{gray95}{gray}{0.95}

\tikzmath{
  function xparabola(\y, \xv, \yv, \S) {
    return 0.5 * ( (\xv + \S) - (\y - \yv) * (\y - \yv) / (\S - \xv) );
  };
  function pairint(\xa, \ya, \xb, \yb, \S) {
    \a = \xb - \xa;
    \b = \yb * (\S - \xa) - \ya * (\S - \xb);
    \c = \yb * \yb * (\S - \xa) - \ya * \ya * (\S - \xb) - (\S - \xa) * (\S - \xb) * (\xb - \xa);
    \d = sqrt( \b * \b - \a * \c);
    \yint = (\b + \d) / \a;
    if \yint < \ya then {
      \yint = (\b - \d) / \a;
    };
    if \yint > \yb then {
      \yint = (\b - \d) / \a;
    };
    return \yint;
  };
}

\newcommand{\arrow}[5]{
  \pgfmathsetmacro{\rx}{#3 - #1}
  \pgfmathsetmacro{\ry}{#4 - #2}
  \pgfmathsetmacro{\n}{sqrt( \rx * \rx + \ry * \ry )}
  
  \pgfmathsetmacro{\vx}{\rx / \n}
  \pgfmathsetmacro{\vy}{\ry / \n}

  \pgfmathsetmacro{\wx}{\vy}
  \pgfmathsetmacro{\wy}{-\vx}  

  \pgfmathsetmacro{\g}{0.025}
  \pgfmathsetmacro{\gx}{-\g * \wx}
  \pgfmathsetmacro{\gy}{-\g * \wy}  
  
  \draw[#5, line width=0.75] 
       (#1 + \gx + 0.75 * \g * \vx, #2 + \gy + 0.75 * \g * \vy) 
    -- (#3 + \gx - 0.75 * \g * \vx, #4 + \gy - 0.75 * \g * \vy); 
    
  \pgfmathsetmacro{\p}{30}
  \pgfmathsetmacro{\d}{0.25}
  \pgfmathsetmacro{\e}{0.1}
  \pgfmathsetmacro{\lx}{#3 - (\d + \e * tan(\p)) * \vx - \e * \wx}
  \pgfmathsetmacro{\ly}{#4 - (\d + \e * tan(\p)) * \vy - \e * \wy}  
  \pgfmathsetmacro{\mx}{#3 - \d * \vx}
  \pgfmathsetmacro{\my}{#4 - \d * \vy}  
  
  \fill[#5]    (#3 + \gx - 0.75 * \g * \vx, #4 + \gy - 0.75 * \g * \vy) 
            -- (\lx + \gx, \ly + \gy)
            -- (\mx + \gx, \my + \gy) 
            -- cycle;
}

\begin{document}

\begin{tikzpicture}

  
  \draw[white] (-2.5, -2.5) rectangle (2.5, 2.5);

  \pgfmathsetmacro{\xa}{1}
  \pgfmathsetmacro{\ya}{0.5}
    
  \pgfmathsetmacro{\xb}{-1.5}
  \pgfmathsetmacro{\yb}{-0.5}
      
  \fill[dark] (\xa, \ya) circle (0.06) node[left] { $s_{n_{2}}$ };
  \fill[dark] (\xb, \yb) circle (0.06) node[left] { $s_{n_{1}}$ };

  \pgfmathsetmacro{\Sf}{1.35}
  \draw[midteal, line width=1] (\Sf, -2.5) -- (\Sf, +2.5);
  \node[midteal] at (\Sf + 0.25, -1.75) { $S$ };
  
  \begin{scope}
    \clip (-2.5, -2.5) rectangle (2.5, 2.5);
    \pgfmathsetmacro{\yinta}{pairint(\xb, \yb, \xa, \ya, \Sf)}
    \pgfmathsetmacro{\yintb}{pairint(\xa, \ya, \xb, \yb, \Sf)}
    
    \draw[domain={-2.5:2.5}, smooth, samples=50, line width=1, variable=\y, color=gray90] 
      plot ( {xparabola(\y, \xa, \ya, \Sf)}, \y );
    
    \draw[domain={-2.5:2.5}, smooth, samples=50, line width=1, variable=\y, color=gray90] 
      plot ( {xparabola(\y, \xb, \yb, \Sf)}, \y );
   
    \draw[domain={-2.5:\yinta}, smooth, samples=50, line width=1, variable=\y, color=mid] 
      plot ( {xparabola(\y, \xb, \yb, \Sf)}, \y );
      
    \draw[domain={\yinta:\yintb}, smooth, samples=50, line width=1, variable=\y, color=mid] 
      plot ( {xparabola(\y, \xa, \ya, \Sf)}, \y );  
      
    \draw[domain={\yintb:2.5}, smooth, samples=50, line width=1, variable=\y, color=mid] 
      plot ( {xparabola(\y, \xb, \yb, \Sf)}, \y );
      
    \pgfmathsetmacro{\xinta}{xparabola(\yinta, \xb, \yb, \Sf)}
    \pgfmathsetmacro{\xintb}{xparabola(\yintb, \xb, \yb, \Sf)}
      
    \fill[dark] (\xinta, \yinta) circle (0.06) node[below right] { $b_{n_{1} n_{2}}$ };
    \fill[dark] (\xintb, \yintb) circle (0.06) node[above right] { $b_{n_{2} n_{1}}$ };
  \end{scope}

\end{tikzpicture}

\end{document}  